../cictikz/src/cictikz/data/tex/cictikz_preamble.tex

% The analog chain: a sensor, an analog front-end, an ADC, bits.  The sensor
% box carries a current source and a voltage source side by side, which is
% the point the prose makes about sensors having their signal in different
% domains.

\begin{circuitikz}[american, thick, transform shape]

  \def\yb{2.05}   % the signal path through the boxes
  \def\ytop{3.2}
  \def\ybot{0.95}

  % ---- boxes ----
  \draw (0.5,\ybot) rectangle (3.8,\ytop);
  \draw (4.75,\ybot) rectangle (8.05,\ytop);
  \draw (9.0,\ybot) rectangle (12.3,\ytop);

  \node[anchor=south] at (2.15,3.32) {Sensor};
  \node[anchor=south] at (6.4,3.32) {AFE};
  \node[anchor=south] at (10.65,3.32) {ADC};

  % ---- what a sensor can look like ----
  % a current source
  \draw (1.52,2.7) -- (1.52,2.36);
  \draw (1.52,2.05) circle (0.31);
  \draw[->] (1.52,2.28) -- (1.52,1.82);
  \draw (1.52,1.74) -- (1.52,1.3);
  % and a voltage source
  \draw (2.46,2.7) -- (2.46,2.37);
  \draw (2.46,2.06) circle (0.31);
  \node at (2.46,2.21) {\scriptsize $+$};
  \node at (2.46,1.92) {\scriptsize $-$};
  \draw (2.46,1.75) -- (2.46,1.3);

  % ---- what the front-end does ----
  \node at (6.4,2.62) {\small Amplification};
  \node at (6.4,2.07) {\small Frequency selectivity};
  \node at (6.4,1.52) {\small Domain transfer};

  % ---- and what the converter does ----
  \node at (10.65,2.55) {\small discrete time?};
  \node at (10.65,1.85) {\small Discrete value};
  \node at (10.65,1.50) {\small (Quantization)};

  % ---- the signal path ----
  \draw[->] (3.8,\yb) -- (4.75,\yb);
  \draw[->] (8.05,\yb) -- (9.0,\yb);
  \draw[->] (12.3,\yb) -- (13.3,\yb);
  \draw (12.78,1.82) -- (13.02,2.32);
  \node[anchor=south west] at (13.1,2.15) {Bits};

\end{circuitikz}
\end{document}
